\documentclass[11pt]{article}

\usepackage[utf8]{inputenc}
\usepackage[T1]{fontenc}
\usepackage{lmodern}
\usepackage{geometry}
\usepackage{amsmath,amssymb,amsfonts,amsthm,mathtools}
\usepackage{enumitem}
\usepackage{microtype}
\usepackage{hyperref}

\geometry{margin=1in}
\hypersetup{colorlinks=true, linkcolor=black, urlcolor=blue, citecolor=black}
\setlist[enumerate]{topsep=0.35em,itemsep=0.3em}
\setlist[itemize]{topsep=0.3em,itemsep=0.25em}

\newcommand{\Aether}{\AE{}ther}
\newcommand{\Aflow}{\AE{}ther-flow}
\newcommand{\TheoryProgram}{\Aether{} / \Aflow{} framework}
\newcommand{\Sub}{\mathcal{A}}
\newcommand{\Step}{\mathsf{Step}}
\newcommand{\RespAlg}{\mathcal{R}_{\Gamma}}
\newcommand{\Rank}{\mathsf{rank}_{\mathcal{S}}}
\newcommand{\Cycle}{\mathcal{C}_{\Gamma}}
\newcommand{\Lang}{\mathcal{L}_{0}}
\newcommand{\Sel}{\mathfrak{S}}
\newcommand{\Dom}{\mathsf{Dom}}
\newcommand{\Inv}{\mathsf{Inv}}
\newcommand{\Fail}{\mathsf{Fail}}

\newtheorem{auditfinding}{Audit finding}
\newtheorem{smugglingrisk}{Smuggling risk}
\newtheorem{promotionblocker}{Promotion blocker}
\newtheorem{auditverdict}{Audit verdict}
\newtheorem{nextburden}{Next burden}

\title{\textbf{Localization Source-Basis Axiom Selector-Domain Smuggling Audit}\\[0.35em]
\large Target-Import Review of Survey Domains and Invariance Criteria for \(D_0\)--\(D_F\)}
\author{Smuggling Auditor AgentJob AJ-RT-20260608-017-001}
\date{June 9, 2026}

\begin{document}

\maketitle

\begin{abstract}
This draft audits the selector-domain repair packet
\((\Dom_0,\ldots,\Dom_4,\Dom_F)\) and its invariance criteria for hidden
target imports. The repair packet is a useful control improvement because it
turns previous underselection findings into explicit survey-domain
obligations and blocks use of family-valued, underdetermined,
benchmark-sensitive, or uncomputed outputs. The audit nevertheless finds that
the repair layer is not promotion-ready. Source-named domain membership,
bounded locality, support equivalence, discriminator closure, quotient
admissibility, transport equivalence, and status-transition rules remain
possible target-import channels unless fixed or surveyed by source-side
criteria. Candidate reconstruction and Gate Chair review remain blocked.
\end{abstract}

\section{Scope and Audit Standard}

This artifact is a draft/control output for task \texttt{RT-20260608-017}.
It does not edit canonical ontology TeX, adopt A0--A4 or \(S_0\)--\(S_F\) as
canonical ontology, authorize candidate reconstruction, issue a Gate Chair
verdict, or reject the broader \TheoryProgram{}.

The audit asks two questions:
\begin{enumerate}
    \item Does the repair packet directly introduce target metric, target
    null-cone, target coordinate, target observer protocol, target rank, or
    target proper-time data as source input?
    \item Even if it avoids direct import, can the declared survey domains,
    invariance criteria, or status transitions be chosen or interpreted so
    that target structures enter before benchmark comparison?
\end{enumerate}

The first question is a declaration screen. The second question is the
promotion-readiness screen. A repair domain can be named in source vocabulary
and still smuggle target content if the admissible class is tuned to recover
target-like locality, rank, or clock behavior.

\section{Positive Transparency Screen}

\begin{auditfinding}[Survey-domain burden is explicit]
The repair packet materially improves the control state. It requires each
selector to declare a survey domain \(\Dom_i\), an invariance criterion
\(\Inv_i\), and a failure status before any later candidate may use the
output. It also states that non-invariant outputs remain negative control
results rather than material for silent repair.
\end{auditfinding}

\begin{auditfinding}[No direct target primitive]
The packet does not explicitly install a Lorentzian metric, continuum
manifold topology, target null cone, radar-coordinate construction, target
proper-time normalization, or fixed rank four as a source primitive.
\end{auditfinding}

These are control successes. They do not prove that the survey domains are
source-forced, benchmark-independent, or safe for candidate reconstruction.

\section{\texorpdfstring{\(D_0\)}{D0}: Presentation-Domain Risks}

\begin{smugglingrisk}[Bounded source accessibility can hide target locality]
\(\Dom_0\) surveys seeds, closure operators, and refinement policies. The
packet requires those objects to be named by source records and declared
before benchmark comparison. The remaining risk is the definition of
``bounded source-accessibility'' and the permitted closure/refinement class.
If the bound, closure operator, or refinement policy is chosen because it
produces target-like local arenas, target locality has entered through domain
membership rather than through an explicit metric.
\end{smugglingrisk}

\begin{promotionblocker}[Presentation-domain blocker]
\(\Dom_0\) is not promotion-ready until the admissible seed, closure, and
refinement class is fixed by source dynamics or until the full declared class
is surveyed and non-isomorphic boundaries, fronts, or response-rank inputs are
preserved as negative results.
\end{promotionblocker}

\section{\texorpdfstring{\(D_1\)}{D1}: Support-Equivalence Risks}

\begin{smugglingrisk}[Support equivalence can encode topology or protocol]
\(\Dom_1\) surveys support-equivalence relations and minimality tests. The
packet forbids target topology, geodesic adjacency, target observer protocols,
and benchmark success. That is necessary. It is not sufficient unless the
support-equivalence class is itself source-fixed. Coarser or finer support
classes can create different transition fronts while remaining source-named.
\end{smugglingrisk}

\begin{promotionblocker}[Transition-domain blocker]
The transition selector cannot support candidate reconstruction until every
admissible support relation and minimality test gives the same
localization-relevant front data, or until a later candidate is formulated to
survive the full family of front data.
\end{promotionblocker}

\section{\texorpdfstring{\(D_2\)}{D2}: Perturbation and Discriminator Risks}

\begin{smugglingrisk}[Closure rules can tune first-response ranks]
\(\Dom_2\) closes over source-local perturbations and discriminators generated
by recurrence words, transition windows, chain-internal response distinctions,
and edit rules. Thresholds and windows must be declared before benchmark
comparison. The remaining risk is that the closure generator itself can be
made broad, narrow, or threshold-sensitive enough to recover favorable
first-response ranks.
\end{smugglingrisk}

\begin{promotionblocker}[Discriminator-closure blocker]
First-response ranks cannot be collapsed to a candidate-supporting datum
unless they are stable across the full declared discriminator domain. A rank
spectrum or discriminator-family spectrum remains a result, not a defect to
remove.
\end{promotionblocker}

\section{\texorpdfstring{\(D_3\)}{D3}: Quotient and Rank Risks}

\begin{smugglingrisk}[Quotient admissibility can manufacture rank stability]
\(\Dom_3\) surveys response-word ground sets, recurrence equivalences,
source relabeling symmetries, and recombination quotients. The packet says no
member may be selected because it yields a target-favorable rank. The
remaining risk is subtler: the admissible quotient class can be defined so
that unfavorable quotient rules are excluded before the rank survey begins.
\end{smugglingrisk}

\begin{promotionblocker}[Rank-domain blocker]
A singleton rank spectrum is usable only if the quotient domain was fixed
without target information and includes all source-admissible quotient rules.
Otherwise apparent rank invariance is a domain-selection result, not
source-side evidence.
\end{promotionblocker}

\section{\texorpdfstring{\(D_4\)}{D4}: Clock-Transport Risks}

\begin{smugglingrisk}[Transport equivalence can encode synchronization]
\(\Dom_4\) surveys shared-support relations, stable-cycle equivalences, and
transport maps. The packet properly excludes target synchronization and
proper-time normalization. The unresolved risk is that exact-support,
declared-overlap, and refinement-compatible transport rules can still encode
a connection-like convention, observer synchronization, or clock scale while
remaining source-named.
\end{smugglingrisk}

\begin{promotionblocker}[Clock-domain blocker]
Transported cycle ratios and holonomy reports cannot support candidate
reconstruction until the transport domain is source-fixed or until the full
declared transport family is surveyed and non-invariant ratios remain
negative control outputs.
\end{promotionblocker}

\section{\texorpdfstring{\(D_F\)}{DF}: Status-Ledger Risks}

\begin{smugglingrisk}[Status transitions can become silent promotion]
\(\Dom_F\) requires metadata for surveyed domains, source inputs, target
inputs excluded, benchmark-comparison time, output family, invariance
criterion, and status. This is the right control surface. The remaining risk
is semantic drift: \(\mathsf{not\_computed}\),
\(\mathsf{family\_valued}\), or \(\mathsf{underdetermined}\) can be treated
as partial success, or \(\mathsf{invariant}\) can be asserted before a public
domain survey exists.
\end{smugglingrisk}

\begin{promotionblocker}[Status-ledger blocker]
An \(\mathsf{invariant}\) status is not evidence unless the surveyed domain,
excluded target inputs, comparison time, and transition rule are recorded in a
reproducible ledger. Status labels without survey records are bookkeeping,
not promotion support.
\end{promotionblocker}

\section{Audit Verdict}

\begin{auditverdict}[Transparent repair, unresolved target-import channels]
The selector-domain repair packet passes a narrow declaration audit: it
exposes the survey-domain burden and does not directly install target metric,
target null-cone, target coordinates, target rank, or target proper-time
normalization. It fails promotion-readiness. The unresolved hidden
target-import channels are:
\begin{enumerate}
    \item \(D_0\) can encode target-like locality through source-accessibility
    bounds, closure operators, or refinement policies;
    \item \(D_1\) can encode target topology or observer protocol through
    support equivalence and minimality tests;
    \item \(D_2\) can tune first-response ranks through closure generators,
    thresholds, and windows;
    \item \(D_3\) can manufacture rank stability through quotient-domain
    admissibility; and
    \item \(D_4\) can import synchronization or clock normalization through
    transport equivalence.
\end{enumerate}
\end{auditverdict}

This is a local control result. It does not reject the \TheoryProgram{} and
does not reject selector-domain repair as a research object. It blocks
candidate reconstruction from the present repair layer.

\section{Next Burden}

\begin{nextburden}[Refuter benchmark-independence test]
The logical next role is a Refuter AgentJob that tests whether the repaired
domains and invariance criteria are benchmark-independent over their
admissible source-defined classes. The test should require explicit variation
over seeds, closures, refinements, support equivalences, discriminator
closures, quotient rules, transport maps, and status-transition ledgers. If
any localization-relevant output is non-invariant, the result should remain a
negative control result or route to another repair packet, not to candidate
reconstruction.
\end{nextburden}

Candidate Constructor, canonical ontology edit, and Gate Chair review remain
premature.

\section*{References}

Ricciardi, A. S. (2026, June 9). \emph{Localization source-basis axiom
selector-domain repair packet: Survey domains, invariance criteria, and
negative-result status rules for S0--SF} [Unpublished manuscript]. The
Aether-Flow Interpretation of Relativity Research Project,
\texttt{research\_control/tasks/RT-20260608-016/artifacts/15\_LOCALIZATION\_SOURCE\_BASIS\_AXIOM\_SELECTOR\_DOMAIN\_REPAIR\_PACKET.tex}.

Ricciardi, A. S. (2026, June 9). \emph{Localization source-basis axiom
selection-law benchmark-independence refutation: Invariance test for
source-internal selector rules} [Unpublished manuscript]. The Aether-Flow
Interpretation of Relativity Research Project,
\texttt{research\_control/tasks/RT-20260608-015/artifacts/14\_LOCALIZATION\_SOURCE\_BASIS\_AXIOM\_SELECTION\_LAW\_BENCHMARK\_INDEPENDENCE\_REFUTATION.tex}.

Ricciardi, A. S. (2026, June 9). \emph{Localization source-basis axiom
selection-law smuggling audit: Target-import review of source-internal
selector rules} [Unpublished manuscript]. The Aether-Flow Interpretation of
Relativity Research Project,
\texttt{research\_control/tasks/RT-20260608-014/artifacts/13\_LOCALIZATION\_SOURCE\_BASIS\_AXIOM\_SELECTION\_LAW\_SMUGGLING\_AUDIT.tex}.

\end{document}
